\chapter{Statistical Problem and Analytic Baselines}
\label{chap:baselines}

This chapter fixes the estimand, effect estimate, standard error, and two
analytic baselines. All three methods evaluated in the thesis use the same
bias-corrected MI difference and the same first-order standard error. They
differ only in the reference distribution used to convert the standardized
statistic into a p-value and confidence interval.

\section{Data, estimand, and assumptions}

Let
\begin{equation}
  N^P=(N^P_{ij}),
  \qquad
  N^Q=(N^Q_{ij})
\end{equation}
be independent \(r\times c\) count tables with totals \(n_P\) and \(n_Q\).
The row and column categories have the same definitions in both populations.
The empirical probabilities are
\begin{equation}
  \widehat p_{ij}=\frac{N^P_{ij}}{n_P},
  \qquad
  \widehat q_{ij}=\frac{N^Q_{ij}}{n_Q}.
\end{equation}

The population parameter is the signed difference
\begin{equation}
  \Delta=I(P)-I(Q).
  \label{eq:delta-population}
\end{equation}
The principal test is two-sided:
\begin{equation}
  H_0:\Delta=0
  \quad\text{against}\quad
  H_1:\Delta\ne0.
  \label{eq:two-sample-null}
\end{equation}

The first-order construction uses the following assumptions.

\begin{assumption}[Sampling]
Observations are independent and identically distributed within each
population, and the two population samples are independent.
\end{assumption}

\begin{assumption}[Aligned fixed alphabets]
The row and column categories are fixed in advance and aligned between the two
tables.
\end{assumption}

\begin{assumption}[Positive support]
Every population cell has positive probability. Empirical zero cells are
allowed, but widespread observed support loss is recorded as a diagnostic.
\end{assumption}

\begin{assumption}[Regular MI]
The first-order influence variances \(V(P)\) and \(V(Q)\) are positive. This
excludes exact independence and other degenerate cases from the regular
Student approximation.
\end{assumption}

All logarithms are natural. MI differences and standard errors are therefore
reported in nats. Changing the logarithm base scales both by the same positive
constant and leaves the standardized statistic unchanged.

\section{Plug-in and bias-corrected MI}

For the first table, define empirical pointwise MI by
\begin{equation}
  \widehat\pmi_P(i,j)
  =\log\left(
  \frac{\widehat p_{ij}}
       {\widehat p_{i+}\widehat p_{+j}}
  \right)
  \qquad\text{when }\widehat p_{ij}>0.
  \label{eq:empirical-pmi}
\end{equation}
The plug-in MI estimate is
\begin{equation}
  \Ihat(P)
  =\sum_{i,j:\widehat p_{ij}>0}
  \widehat p_{ij}\widehat\pmi_P(i,j).
  \label{eq:plugin-mi}
\end{equation}
The same calculation gives \(\Ihat(Q)\).

Let
\begin{equation}
  d=(r-1)(c-1).
\end{equation}
Under fixed positive support, the leading bias correction gives
\begin{equation}
  \widehat I_{\mathrm{BC}}(P)
  =\Ihat(P)-\frac{d}{2n_P},
  \qquad
  \widehat I_{\mathrm{BC}}(Q)
  =\Ihat(Q)-\frac{d}{2n_Q}.
  \label{eq:bias-corrected-mi}
\end{equation}
The estimated effect is
\begin{equation}
  \Deltahat
  =\widehat I_{\mathrm{BC}}(P)
  -\widehat I_{\mathrm{BC}}(Q).
  \label{eq:bias-corrected-difference}
\end{equation}

The correction uses the declared alphabet dimensions rather than the observed
nonempty dimensions. It may produce a negative corrected MI estimate in a
sparse sample. The implementation does not clip that estimate at zero because
clipping would alter both the estimand and the sampling distribution.

\section{First-order variance of plug-in MI}

For population \(P\), define
\begin{equation}
  V(P)
  =\sum_{i,j}p_{ij}
  \{\pmi_P(i,j)-I(P)\}^2.
  \label{eq:population-mi-variance}
\end{equation}
This is the population variance of pointwise MI. It governs the sampling
variance of the complete plug-in estimator because the first-order MI error is
an average of centred pointwise-MI contributions.

To see this directly, write one observation as
\(Z_a^{(P)}=(X_a^{(P)},Y_a^{(P)})\). A first-order expansion gives
\begin{equation}
  \Ihat(P)-I(P)
  =
  \frac{1}{n_P}
  \sum_{a=1}^{n_P}
  \{\pmi_P(Z_a^{(P)})-I(P)\}
  +o_p(n_P^{-1/2}).
  \label{eq:mi-asymptotic-linear}
\end{equation}
The summands are independent, have mean zero, and have variance \(V(P)\).
Therefore,
\begin{align}
  \Var\{\Ihat(P)\}
  &\approx
  \Var\left[
  \frac{1}{n_P}
  \sum_{a=1}^{n_P}
  \{\pmi_P(Z_a^{(P)})-I(P)\}
  \right]\\
  &=\frac{1}{n_P^2}
  \sum_{a=1}^{n_P}V(P)\\
  &=\frac{V(P)}{n_P}.
  \label{eq:mi-sampling-variance}
\end{align}
Subtracting the fixed leading bias term in Equation
\eqref{eq:bias-corrected-mi} does not change this first-order variance.

The observed table estimates \(V(P)\) by
\begin{equation}
  \Vhat(P)
  =\sum_{i,j:\widehat p_{ij}>0}
  \widehat p_{ij}
  \{\widehat\pmi_P(i,j)-\Ihat(P)\}^2.
  \label{eq:estimated-mi-variance}
\end{equation}
Define \(\Vhat(Q)\) analogously.

\section{Shared standard error and test statistic}

Because the two samples are independent, their estimated variance
contributions add:
\begin{equation}
  \SEhat^2
  =\frac{\Vhat(P)}{n_P}
  +\frac{\Vhat(Q)}{n_Q}.
  \label{eq:shared-se}
\end{equation}
The common standardized statistic is
\begin{equation}
  T=\frac{\Deltahat}{\SEhat}.
  \label{eq:shared-t}
\end{equation}
Equations~\eqref{eq:bias-corrected-difference}--\eqref{eq:shared-t} are shared
by normal Wald, simple Welch, and expanded Welch. This controlled comparison
isolates the effect of the reference distribution.

\section{Baseline 1: normal Wald}

Under the regular first-order asymptotics,
\begin{equation}
  T\xrightarrow{d}N(0,1)
  \qquad\text{under }H_0.
\end{equation}
The two-sided normal Wald p-value is
\begin{equation}
  p_{\mathrm{Wald}}
  =2\{1-\Phi(|T|)\},
  \label{eq:wald-pvalue}
\end{equation}
where \(\Phi\) is the standard normal distribution function. A
\(100(1-\alpha)\%\) confidence interval for \(\Delta\) is
\begin{equation}
  \Deltahat
  \ \pm\
  z_{1-\alpha/2}\SEhat.
  \label{eq:wald-ci}
\end{equation}

The normal reference is the large-sample limit shared by all three analytic
methods. It does not assign a separate finite-sample uncertainty to
\(\Vhat(P)\) or \(\Vhat(Q)\).

\section{Baseline 2: simple Welch--Satterthwaite}

Simple Welch treats each empirical MI influence variance as though it had the
same component degrees of freedom as an ordinary sample variance:
\begin{equation}
  \nu_P^{\mathrm{simple}}=n_P-1,
  \qquad
  \nu_Q^{\mathrm{simple}}=n_Q-1.
\end{equation}
Substitution into the Welch--Satterthwaite equation gives
\begin{equation}
  \widehat\nu_{\mathrm{simple}}
  =
  \frac{
  \left\{\Vhat(P)/n_P+\Vhat(Q)/n_Q\right\}^2
  }{
  \left\{\Vhat(P)/n_P\right\}^2/(n_P-1)
  +
  \left\{\Vhat(Q)/n_Q\right\}^2/(n_Q-1)
  }.
  \label{eq:simple-welch-df}
\end{equation}
The two-sided p-value is
\begin{equation}
  p_{\mathrm{simple}}
  =2\{1-F_{t_{\widehat\nu_{\mathrm{simple}}}}(|T|)\}.
  \label{eq:simple-welch-p}
\end{equation}

The Student reference has heavier tails than the normal reference, so simple
Welch is weakly more conservative at a fixed observed \(|T|\). Its limitation
is the component degrees of freedom. The value \(n-1\) has a classical basis
for an ordinary sample variance under normal sampling. The quantity
\(\Vhat(P)\) is instead a nonlinear plug-in functional whose pointwise scores
change with the empirical table.

\section{What the expanded method must calculate}

The shared statistic already supplies the estimated effect and its first-order
standard error. The remaining task is to replace the generic values
\(n_P-1\) and \(n_Q-1\) by component degrees of freedom that describe the
sampling reliability of \(\Vhat(P)\) and \(\Vhat(Q)\).

Satterthwaite moment matching requires the first two moments of each positive
variance estimate. To first order,
\begin{equation}
  \E\{\Vhat(P)\}\approx V(P).
\end{equation}
Chapter~\ref{chap:expanded} derives
\(\Var\{\Vhat(P)\}\) by calculating how the complete variance functional
changes when one observation changes the joint table.
